%! Parent Functions and Transformations — Unit 2, Lesson 2.2 — Exit Ticket (Answer Key)
\documentclass[10pt]{article}
\usepackage{atda-article}
\usepackage{atda-key}   % pulls in -boxes; do NOT also load -boxes
\newcommand{\TallMath}[1]{$\displaystyle #1\rule[-1.4em]{0pt}{3.2em}$}
\setlength{\workrowsep}{8pt}

\begin{document}

\pageheader{Unit 2, Lesson 2.2}{Exit Ticket --- Answer Key}
% NAMESTRIP: no \namedateperiod here — mirrors the blank. Teacher-only prose goes
% in the lesson plan as \begin{teachernote}[Exit Ticket], never in a key.

\vspace{0.15in}

\begin{notesbox}{Name the Parent, Name the Move}
\small
Notes closed. Three items.

\begin{center}
\begin{tikzpicture}
\begin{axis}[
    width=9.8cm, height=4.4cm,
    xlabel={$x$}, ylabel={$y$},
    xmin=-3.5, xmax=6.5, ymin=0, ymax=9.6,
    xtick={-3,-2,-1,0,1,2,3,4,5,6}, ytick={0,3,6,9},
    grid=both, grid style={linegray!45}, axis lines=left,
    tick label style={font=\scriptsize}, label style={font=\scriptsize},
    legend entries={{$y=f(x)$ (the parent)},{$y=g(x)$}},
    legend style={font=\scriptsize, draw=none, fill=none,
      at={(0.5,1.02)}, anchor=south, legend columns=2}]
  \addplot[royal, very thick, domain=-3:3, samples=80]{x^2};
  \addplot[goldacc, very thick, dashed, domain=0:6, samples=80]{(x-3)^2};
\end{axis}
\end{tikzpicture}
\end{center}

\begin{enumerate}[leftmargin=1.6em, itemsep=8pt, topsep=4pt]

\item Name the parent function $f$ drawn above. Then describe how $g$ was made from it --- which
way and how far --- and write $g$ two ways: in function notation starting from $f$, and as an
equation in $x$.

\ansline{The quadratic parent, $f(x)=x^2$. Its graph was translated \textbf{right $3$}.}\\
\ansline{$g(x)=f(x-3)$, that is $g(x)=(x-3)^2$. Turning point $(3,0)$.}\\

\item No graph for this one: $y=2^x-5$. Name the parent, describe the transformation, and state
which of the \textbf{domain} and the \textbf{range} moved. Give both in interval notation.

\ansline{Parent $h(x)=2^x$, translated \textbf{down $5$} (the $-5$ is outside the parentheses).}\\
\ansline{Domain $(-\infty,\infty)$ --- unchanged.}\\
\ansline{Range moved from $(0,\infty)$ to $(-5,\infty)$.}\\

\item A classmate says: ``The equation for $g$ has a minus sign in it, so the graph should have
moved \emph{left}.'' They are wrong. Using an input of your own choosing, explain what the minus
sign actually does.

\ansline{The $-3$ is subtracted from the input \emph{before} squaring. At $x=5$, $g$ computes}\\
\ansline{$f(5-3)=f(2)=4$ --- so the new graph at $5$ shows what the parent showed at $2$.}\\
\ansline{Every value has been copied $3$ units farther right, so the graph moved \textbf{right}.}\\

\end{enumerate}
\end{notesbox}

\end{document}
